\section{Related Work}
\label{sec:related-work}

\begin{table*}[t]
	\centering
	\caption{Positioning of representative approaches against policy-private
		consent enforcement under public-chain observability. ``Yes'' indicates explicit
		support, ``Partial'' indicates indirect or incomplete support, and ``No''
		indicates that the dimension is not addressed---typically because it lies outside
		the approach's original design goals; assumptions in parentheses are part of the
		claim.}
	\label{tab:related-work-gap} \footnotesize
	\begin{tabular}{@{}>{\raggedright\arraybackslash}p{0.2\linewidth}>{\raggedright\arraybackslash}p{0.1\linewidth}>{\raggedright\arraybackslash}p{0.1\linewidth}>{\raggedright\arraybackslash}p{0.1\linewidth}p{0.4\linewidth}@{}}
		\hline
		\textbf{Approach}                                                                                                                        & \textbf{Policy confidentiality} & \textbf{Metadata-leakage mitigation} & \textbf{Verifiable auditability} & \textbf{Main limitation}                                                                                                                  \\
		\hline
		Blockchain consent management~\cite{kakarlapudi2021systematic,welzel2025enabling}                                                        & No                              & No                                   & Yes                              & Records and enforces consent verifiably on chain, but policy confidentiality and policy-dependent leakage are typically out of scope.     \\
		Data / transaction privacy~\cite{bethencourt2007ciphertext,ateniese2006improved,sasson2014zerocash,rondelet2019zeth,williamson2018aztec} & No                              & Partial                              & Partial                          & Protects encrypted data access or transaction semantics, but not the governing consent policy or policy-dependent authorization workflow. \\
		Hidden-policy access control~\cite{wu2023blockchain,hu2023access}                                                                        & Partial                         & Partial                              & Partial                          & Reduces direct policy exposure, but may leak through workflows, updates, or policy-holder interactions.                                   \\
		MPC / HE policy evaluation~\cite{ardagna2009xacml,sheikhalishahi2019privacy}                                                             & Yes                             & No                                   & Partial                          & Supports confidential computation, but usually lacks public-chain-verifiable authorization evidence.                                      \\
		Universal SNARKs / zkVMs~\cite{gabizon2019plonk,campanelli2019legosnark,liu2025ceno}                                                     & No                              & Partial                              & Yes                              & Supports proof reuse and public verification, but does not by itself hide policy-dependent public artifacts.                              \\
		Metadata-leakage defenses~\cite{miller2023vtbc}                                                                                          & No                              & Yes                                  & No                               & Models or mitigates observable signals, but does not provide hidden-policy consent enforcement.                                           \\
		TPS-ZKPI                                                                                                                                 & Yes (under TPS trust)           & Yes (modeled channels)               & Yes                              & Provides policy-private, metadata-aware, verifiable consent enforcement under the stated TPS trust assumption.                            \\
		\hline
	\end{tabular}
\end{table*}

Blockchain-based consent management for healthcare data sharing has been studied extensively, with smart contracts used to record, enforce, and audit consent across institutions~\cite{kakarlapudi2021systematic, welzel2025enabling,li2024you}.
This line of work focuses on making consent decisions verifiable and auditable, and the privacy mechanisms surveyed in this space mainly protect health data contents and user identities~\cite{li2025privacy}.
The confidentiality of the governing consent policy---and what public-chain enforcement artifacts reveal about it---is typically out of scope.

Beyond the consent-management setting, existing work addresses important fragments of policy-private consent enforcement, but not their combination under public-chain observability.
Encrypted data-sharing mechanisms, such as Ciphertext-Policy Attribute-Based Encryption (CP-ABE) and proxy re-encryption (PRE)~\cite{bethencourt2007ciphertext,ateniese2006improved}, control access to encrypted data. However, standard CP-ABE attaches the access policy to the ciphertext in the clear, exposing the policy itself, while PRE-based delegation may reveal authorization patterns that data encryption is not designed to hide.
Transaction-privacy systems, including Zerocash-style protocols~\cite{sasson2014zerocash} and Ethereum-oriented systems such as ZETH and AZTEC~\cite{rondelet2019zeth,williamson2018aztec}, provide public verification over private transaction witnesses, but they mainly target asset ownership, values, or participant relationships.

Hidden-policy access-control work reduces direct policy exposure by using encrypted policies, hidden-policy ABE variants, off-chain matching, or policy-holder components~\cite{wu2023blockchain,hu2023access}.
Functional encryption gives a more general cryptographic abstraction for revealing function
outputs while hiding inputs~\cite{boneh2011functional}, but it does not by itself yield publicly verifiable authorization evidence.
Private-state blockchain systems, including Secret Network and Oasis/Ekiden, protect smart-contract state or execution through encrypted state and trusted execution environments~\cite{secretnetwork_docs, cheng2019ekiden}.
These lines of work address adjacent forms of policy or state hiding, but TPS-ZKPI
targets a more specific setting: hiding consent policies on a public chain while
retaining public verifiability of each authorization decision.
Even where policies are hidden, existing designs may still leak policy-dependent artifacts through ciphertext structures, policy-update transactions, authorization workflows, or repeated interactions.
Moreover, many target permissioned or off-chain settings, where enforcement artifacts are not globally visible as they are on public blockchains.

Privacy-preserving access-control and policy-evaluation mechanisms have been studied in credential-based and XACML/SAML settings~\cite{ardagna2009xacml}, as well as in multi-party policy evaluation using homomorphic encryption and secure function evaluation~\cite{sheikhalishahi2019privacy}.
These approaches support confidential authorization computation, but they are not primarily designed for settings where policy-evaluation artifacts---calldata, public inputs, emitted events, timing, and gas usage---are permanently and globally visible on a public blockchain.
In addition, these approaches do not directly provide compact evidence that a smart contract can verify against a committed, provider-defined policy instance.

Reusable proving systems, including universal-setup SNARKs, composable proof systems, and zkVMs~\cite{gabizon2019plonk,campanelli2019legosnark,liu2025ceno}, improve maintainability by reducing the need for policy-specific circuits or verifier contracts.
However, when deployed for public-chain consent enforcement, reusable proving alone does not guarantee policy privacy: policy-dependent program representations, public inputs, calldata layouts, logs, or gas profiles may still remain observable.
This motivates TPS-ZKPI's use of a domain-specific reusable policy interpreter, which preserves proof reuse while exposing only a policy-independent public verification interface.

Public-chain metadata-leakage studies show that cryptographic privacy
guarantees may be weakened by observable transaction-level artifacts such as
timing and volume patterns~\cite{beres2021blockchain}. Techniques such as batching,
decoys, and traffic shaping can reduce some observable signals~\cite{miller2023vtbc}, but they are
usually designed for payment privacy, anonymity sets, or transaction-linkability
settings. They do not directly provide hidden-policy compliance verification
or auditable consent enforcement. In contrast, TPS-ZKPI treats metadata leakage as
part of the policy-privacy objective.

Existing approaches address verifiable consent recording, data confidentiality,
transaction privacy, policy hiding, confidential policy evaluation, proof reuse, or
metadata mitigation in isolation,
whereas policy-private consent enforcement under public-chain observability requires
the combination of policy confidentiality, metadata-leakage control, verifiable
auditability, and a stable verification interface for evolving Boolean-based consent
policies. Table~\ref{tab:related-work-gap}
summarizes this positioning.


